\section{Exact finite-dimensional rank-one attainment}\label{ex:sec:finite-attainment}
The closure has a sharp finite-condition-number refinement in rank one.
For $\tau=0$ the relative spectrum is identically one; assume $\tau>0$.

\begin{theorem}[Exact rank-one realization]\label{ex:thm:exact}
Fix $n\ge1$, $0<\alpha<1$, and $\tau>0$.
A decreasing list $q$ is the relative spectrum of
\[
 A^{-\alpha/2}(A+uu^T)^\alpha A^{-\alpha/2},
 \qquad A\succ0,\quad u^TA^{-1}u=\tau,
\]
if and only if for some $1\le m\le n$,
\[
 q_1>\cdots>q_m>1,\qquad
 q_{m+1}=\cdots=q_n=1,\qquad
 \prod_iq_i=(1+\tau)^\alpha.
\]
In particular every repeated value strictly above one occurs only in
the closure. Every admissible list with simple values above one is
attained at finite positive condition number.
\end{theorem}

\begin{proof}
Let $\mathcal K$ be the cyclic subspace of $(A,u)$ and write
$m=\dim\mathcal K$. It and its orthogonal complement reduce both
endpoints. On $\mathcal K^\perp$ the correction $C=Q_\alpha-I$ is zero.
On $\mathcal K$ take an eigenbasis with
$A=\diag(\theta_1<\cdots<\theta_m)$ and $u_i>0$.
The rank-one fractional resolvent identity gives
\[
 C_{ij}=h_i h_j\int_0^\infty
 \frac{w(t)}{(\theta_i+t)(\theta_j+t)}\,dt,
 \quad h_i=u_i\theta_i^{-\alpha/2}>0,
 \quad
 w(t)=\frac{\sin(\pi\alpha)}\pi
 \frac{t^\alpha}{1+\sum_i u_i^2/(\theta_i+t)}>0.
\]
The integral converges and is a Gram matrix. For any row and column
index sets $I,J$ of size $k$, continuous Cauchy--Binet gives
\[
 \frac{\det C[I,J]}{\prod_{i\in I}h_i\prod_{j\in J}h_j}
 =\int_{0<t_1<\cdots<t_k<\infty}
 \det\left[\frac1{\theta_i+t_l}\right]_{i\in I,l\le k}
 \det\left[\frac1{\theta_j+t_l}\right]_{j\in J,l\le k}
 \prod_{l=1}^k w(t_l)\,dt_1\cdots dt_k>0.
\]
The Cauchy determinant formula makes both determinants strictly positive.
One can first truncate to a compact integration interval, then pass
to the limit; the integrand is nonnegative and the entries converge.
Thus the active correction is strictly totally positive, and positive
definite. If its eigenvalues are
$\lambda_1\ge\cdots\ge\lambda_m>0$, every compound matrix has
strictly positive entries. Perron--Frobenius applied to its $k$th
compound implies $\lambda_k>\lambda_{k+1}$ for $k<m$:
otherwise its largest eigenvalue $\lambda_1\cdots\lambda_k$ would be
repeated. This is the classical oscillation-matrix simplicity mechanism~\cite{Pinkus}.
The determinant lemma gives the required product.

For sufficiency first work in the active dimension $m$.
Put $S_0=\log(1+\tau)$ and $x_i^*=\alpha^{-1}\log q_i$.
Choose a small closed Euclidean ball $K$ centered at $x^*$ in
$\mathcal H=\{x:\sum_i x_i=S_0\}$, of radius $\rho>0$,
contained in $x_1>\cdots>x_m>0$. For $x\in K$ set
\[
 z_i(x)=(e^{x_i}-1)e^{\sum_{j>i}x_j},\quad
 A_\delta=\diag(\delta^{m-i}),\quad
 u_i(x)=\sqrt{(A_\delta)_{ii}z_i(x)}.
\]
The telescoping identity gives $\sum_i z_i(x)=e^{S_0}-1=\tau$.
Write $s_i(x)=\sqrt{z_i(x)}$. The reverse pivots of
$I+s(x)s(x)^T$ are $e^{x_i}$.
Lemma~\ref{ah:lem:pivots}, uniformly for this compact family, gives
\[
 \sup_{x\in K}\norm{Q_\delta(x)-\diag(e^{\alpha x_i})}
 \longrightarrow0.
\]
Consequently the continuous ordered-log-spectrum map
\[
 F_\delta(x)=\alpha^{-1}\log\lambda(Q_\delta(x))
\]
takes values in $\mathcal H$ and converges uniformly to $x$ on $K$.
The affine-space condition follows from the exact determinant identity.
For sufficiently small positive $\delta$,
$T_\delta(x)=x^*+x-F_\delta(x)$ maps $K$ to itself, because
$\norm{T_\delta(x)-x^*}_2<\rho$.
Brouwer gives a fixed point, at which $F_\delta(x)=x^*$.
For $m=1$ use the scalar identity directly. Finally pad inactive
dimensions with any SPD block and zeros in $u$.
\end{proof}

The proof is existential and supplies no useful upper bound on the
condition number of an exact witness. The simplicity statement uses
classical Cauchy-kernel total positivity, while the inverse existence
uses Brouwer; historical priority for this application has not been
established.

\begin{corollary}[Exact rank and simplicity for one graph edge]
\label{ex:cor:graphsimple}
For a positive edge-conductance edit of a connected graph, every
nonzero eigenvalue of the relative fractional correction is positive
and simple. Its rank is
\[
 \#\{\lambda>0:\lambda\in\spec(L),\ P_\lambda b\ne0\},
\]
where $b$ is the edited edge incidence and $P_\lambda$ the spectral
projector of $L$. For an endpoint edge of an $N$-vertex weighted path,
the rank is exactly $N-1$ for every positive edit size.
\end{corollary}
\begin{proof}
Restrict to $\mathbf1^\perp$ and apply the cyclic-subspace argument
above. A weighted path Laplacian is irreducible symmetric tridiagonal,
so its eigenvalues are simple. If an eigenvector $v$ with positive
eigenvalue is orthogonal to an endpoint incidence, then $v_1=v_2$.
The endpoint eigenvector equation gives $\lambda v_1=0$, hence
$v_1=v_2=0$. Recurrence along the path forces $v=0$, a contradiction.
Every positive path eigenspace is therefore seen by that incidence.
\end{proof}

This exact graph obstruction is compatible with low approximate rank.
The path closure theorem permits arbitrarily many eigenvalues tending
to zero. It does not assert exact path realization of each matrix
spectrum with exactly zero padded response eigenvalues.

